\documentclass[conference]{IEEEtran}
\IEEEoverridecommandlockouts
% The preceding line is only needed to identify funding in the first footnote. If that is unneeded, please comment it out.
\usepackage{cite}
\usepackage{amsmath,amssymb,amsfonts}
\usepackage{algorithmic}
\usepackage{graphicx}
\usepackage{textcomp}
\usepackage{xcolor}
\def\BibTeX{{\rm B\kern-.05em{\sc i\kern-.025em b}\kern-.08em
    T\kern-.1667em\lower.7ex\hbox{E}\kern-.125emX}}
\begin{document}

\title{CPE380 Project 3: Pipelined Instruction Implementation}

\author{\IEEEauthorblockN{1\textsuperscript{st} Xavier Jones}
\IEEEauthorblockA{\textit{Department of Electrical and Computer Engineering} \\
\textit{University of Kentucky College of Engineering}\\
Lexington, Kentucky, USA \\
xcjo222@uky.edu}
\and
\IEEEauthorblockN{2\textsuperscript{nd} Benjamin Carey}
\IEEEauthorblockA{\textit{Department of Computer Science} \\
\textit{University of Kentucky College of Engineering}\\
Lexington, Kentucky, USA \\
bjca251@uky.edu}
\and
\IEEEauthorblockN{3\textsuperscript{rd} Jiawen Liu}
\IEEEauthorblockA{\textit{Department of Electrical and Computer Engineering} \\
\textit{University of Kentucky }\\
Lexington, Kentucky, USA \\
jli416@uky.edu}
}

\maketitle

\begin{abstract}
    This document outlines the design used to extend a Verilog multi-cycle
    pipelined processor simulator. The objective of this project was to add
    support for integer minimum (\texttt{min}), integer multiplication (\texttt{mult}), and jump
    register (\texttt{jr}) functionalities.
\end{abstract}

\begin{IEEEkeywords}
    
\end{IEEEkeywords}

\section{Introduction}
This assignment was proposed as an introduction to the pipelined processor
design taught in Dr. Hank Dietz's Computer Organization course. Herein, one
learns all the components associated with a high-level understanding of
processor design, advanced digital logic, and pipelining a processor's
execution stages. Our assignment was to make a minimal number of changes to the
instruction set architecture of a pipelined MIPS-based processor to add an
integer minimum and integer multiplication instruction as well as a control
flow instruction that jumps the instruction register to the value in a
register.

\section{Implemented Instructions}
\subsection{The \texttt{min} Instruction}
The \texttt{min} instruction compares the values in two registers and stores
whichever value is lesser in the destination register. It is encoded as an
\textbf{R-type} instruction with a \texttt{FUNCT} field of \texttt{0x01}. The
two modifications to the code were as follows:
\begin{itemize}
    \item Firstly, the \textbf{decode stage} is modified to catch the case when
        the \texttt{FUNCT} field reads an encoded \texttt{MIN} value. As we
        store the result value in a register, we set the write-back flags to
        suit (\texttt{RegDst=1, RegWrite=1}). Furthermore, a new ALU operation
        \texttt{ALUMIN} is defined to signal the execution stage which ALU
        operation to execute.
    \item Secondly, the \textbf{execute stage} is modified to compute the
        minimum value between its two inputs. This is performed using a ternary
        operator \texttt{(ID\_s < ID\_src) ? ID\_s : ID\_src}.
\end{itemize}
Note that a min instruction that finds the minimum between a register and an
immediate value can be trivially implemented with the changes made here.

\subsection{The \texttt{mult} Instruction}
The \texttt{mult} instruction performs an integer multiplication with values in
two registers and stores the result in the destination register. It is encoded
as an \textbf{R-type} instruction with a \texttt{FUNCT} field of \texttt{0x18}.
The two modifications to the code were as follows:
\begin{itemize}
    \item Firstly, the \textbf{decode stage} is modified to catch the case when
        the \texttt{FUNCT} field reads an encoded \texttt{MULT} value. As we
        store the result value in a register, we set the write-back flags to
        suit (\texttt{RegDst=1, RegWrite=1}). Furthermore, a new ALU operation
        \texttt{ALUMULT} is defined to signal the execution stage which ALU
        operation to execute.
    \item Secondly, the \textbf{execute stage} is modified to compute the
        product. This is performed using Verilog's native multiplication
        operator \texttt{ID\_s * ID\_src}.
\end{itemize}
Note that a mult instruction that finds the product between a register and an
immediate value can also be trivially implemented with the changes made here.

\subsection{The \texttt{jr} Instruction}
The \texttt{jr} instruction updates the program counter to the value stored in
the input register \texttt{\$rs}. It is encoded as an \textbf{R-type}
instruction with a \texttt{FUNCT} field of \texttt{0x08}. Implementing this
required the modification of the following phases of the pipelined architecture:
\begin{itemize}
    \item Primarily, the \textbf{decode stage} is modified to catch the case when
        the \texttt{FUNCT} field reads an encoded \texttt{JR} value. The
        control logic of the decode stage was modified to signal a
        \texttt{JumpReg} flag whenever the jump register is called. This is
        used to multiplex the target address between the branch-or-next MUX and
        the value stored in the \texttt{\$rs} register. As the processor's
        branch prediction automatically assumes the branch is not taken, the
        predictor's \texttt{squash} flag is set to true.
    \item All the other stages are left unmodified as no registers or memory
        addresses are operated on or modified.
\end{itemize}

\section{Testing and Verification}
Testing was conducted by modifying the initial begin block of the processor
module to preload specific test cases into memory (the \texttt{m} register
block in the Verilog code). Trace outputs were checked to ensure the
\textbf{write back} (\texttt{WB}) stage properly committed expected values.

\section{Known Issues}
During testing of the \texttt{jr} instruction, there was a structural/data hazard. The
test sequence initially used an \texttt{addiu} instruction to load the target memory
address into a register, immediately followed by the \texttt{jr} instruction attempting
to read that register.

Because the baseline pipeline forwarding logic does not safely handle immediate
R-to-branch data dependencies, the processor fetched an incorrect ALU artifact
instead of the resolved address. This was fixed for the local tests by
introducing independent instructions between the \texttt{addiu} and \texttt{jr}
instructions in the testbench. This allowed the address to safely get to the
register file.

\end{document}
